\cvsection{Projects}
\begin{cventries}
	
	\cventry
	{}
	{Recommendation System for PundA (푼다를 위한 추천 시스템)}
	{Apr. 2019 - Jun. 2019}
	{}
	{
		\begin{cvitemsForEducation}
			\item[] {\textbf{Technologies used: Python, Tensorflow}}
			\item {Developing the problem track's recommendation system (Tensorflow) to improve the PundA product developed by Mathpresso Inc.}
		\end{cvitemsForEducation}
	}
	
	\cventry
    {}
    {Explainable AI (XAI) Project}
    {Mar. 2018 - Des. 2018}
    {}
    {
      \begin{cvitemsForEducation}
        \item[] {\textbf{Technologies used: Python, TkInter, TensorFlow}}
        \item {Developing both the GUI for the Interactive Machine Learning (TkInter) and the algorithm to correct the prediction based on the prediction's explanation provided by LIME Algorithm (TensorFlow)}
      \end{cvitemsForEducation}
    }
	
	\cventry
	{in cooperation with Kookmin University, Seoul, South Korea}
	{ERC Project}
	{Jun. 2017 - Nov. 2017}
	{}
	{
		\begin{cvitemsForEducation}
			\item[] {}
			\item[] {\textbf{Technologies used: Python, ScikitLearn, TensorFlow}}
			\item {Classification towards the Sleep-Stage Pattern using Machine Learning}
			\item {Comparing the Machine Learning Algorithm Performance (Scikit-Learn) with the Deep Learning Algorithm Performance (TensorFlow)}
		\end{cvitemsForEducation}
	}

	\cventry
	{}
	{AI-Doctor Project}
	{Mar. 2017 - Jun. 2017}
	{}
	{
		\begin{cvitemsForEducation}
			\item[] {\textbf{Technologies used: Python, TensorFlow}}
			\item {Using Deep Learning Algorithm (especially RNN and CNN) to classify whether the patient has diseases (for example, Arrhytmia) (TensorFlow)}
		\end{cvitemsForEducation}
	}
	
	\cventry
	{in accordance towards the completion of the CS408 (Computer Science Project) course module}
	{BroCards Project (Team)}
	{Sep. 2016 - Dec. 2016}
	{}
	{
		\begin{cvitemsForEducation}
			\item[] {}
			\item[] {\textbf{Technologies used: Java, Android Studio}}
			\item {Developing the Android Application for the WiFi-based Card Game (mostly using Java and JSon module)}
		\end{cvitemsForEducation}
	}
	

\end{cventries}
